% 第九卷 13.3.5：基于大语言模型与函数调用的任务型对话（跨栏，约160 mm宽）。
% FigureInput：用户；FigureModel：大语言模型；FigureOutput：执行器与服务；虚线框：可选的显式状态记录。
\begin{tikzpicture}[
  x=1mm,y=1mm,
  figure picture,
  user/.style={figure node,draw=FigureInput,fill=FigureInputFill,minimum height=14mm,
    minimum width=16mm,inner xsep=2pt},
  model/.style={figure node,draw=FigureModel,fill=FigureModelFill,minimum height=14mm,
    minimum width=40mm,inner xsep=2pt,line width=.9pt},
  exec/.style={figure node,draw=FigureOutput,fill=FigureOutputFill,minimum height=14mm,
    minimum width=30mm,inner xsep=2pt},
  record/.style={figure node,draw=FigureMuted,fill=white,minimum height=8mm,
    minimum width=56mm,inner xsep=2pt,dash pattern=on 2pt off 1.5pt}
]
  \node[user] (user) at (8,30) {用户};
  \node[model] (llm) at (56,30) {大语言模型\\系统提示＋函数模式＋历史};
  \node[exec] (exe) at (110,30) {执行器\\校验参数与权限\\注入认证身份};
  \node[exec] (svc) at (148,30) {服务\\\texttt{check\_registration}};
  \node[record] (state) at (56,6) {状态记录（可选）：日期、人数及来源轮次};

  \draw[figure flow] ([yshift=3mm]user.east) -- node[figure edge label,above=1pt] {话语} ([yshift=3mm]llm.west);
  \draw[figure flow] ([yshift=-3mm]llm.west) -- node[figure edge label,below=1pt] {回复} ([yshift=-3mm]user.east);
  \draw[figure flow] ([yshift=3mm]llm.east) -- node[figure edge label,above=1pt] {函数调用} ([yshift=3mm]exe.west);
  \draw[figure flow] ([yshift=-3mm]exe.west) -- node[figure edge label,below=1pt] {工具消息} ([yshift=-3mm]llm.east);
  \draw[figure flow] ([yshift=3mm]exe.east) -- node[figure edge label,above=1pt] {调用} ([yshift=3mm]svc.west);
  \draw[figure flow] ([yshift=-3mm]svc.west) -- node[figure edge label,below=1pt] {返回} ([yshift=-3mm]exe.east);
  \draw[figure flow,dash pattern=on 2pt off 1.5pt] (llm.south) -- node[figure note,right=1pt] {写出与核对} (state.north);
\end{tikzpicture}
